\documentclass[11pt]{article}
\usepackage[margin=1in]{geometry}
\usepackage{amsmath,amssymb,amsthm,mathtools}
\usepackage{lmodern}
\usepackage{booktabs}
\usepackage{enumitem}
\usepackage{xcolor}
\usepackage[hidelinks]{hyperref}
\hypersetup{
  pdftitle={Moving-Hole Defects in Translated Reduced-Residue BDH Variances},
  pdfauthor={Liang Wang},
  pdfsubject={Exact moving-hole compiler and deterministic critical block bound}
}

\newtheorem{theorem}{Theorem}[section]
\newtheorem{proposition}[theorem]{Proposition}
\newtheorem{lemma}[theorem]{Lemma}
\newtheorem{corollary}[theorem]{Corollary}
\newtheorem{remark}[theorem]{Remark}
\newtheorem{definition}[theorem]{Definition}

\newcommand{\Fq}{\mathbb Z/q\mathbb Z}
\newcommand{\e}{\mathrm e}
\newcommand{\one}{\mathbf 1}
\newcommand{\cQ}{\mathcal Q}
\newcommand{\cV}{\mathcal V}
\newcommand{\cM}{\mathcal M}
\newcommand{\avg}{\operatorname{Av}}

\title{\textbf{Moving-Hole Defects in Translated Reduced-Residue BDH Variances:}\\
A Rank-Two Formula and Critical Block Bound}
\author{Liang Wang\\
Huazhong University of Science and Technology\\
Wuhan 430074, P.R. China}
\date{August 2026}

\begin{document}
\emergencystretch=2em
\maketitle

\begin{abstract}
Reduced-residue Barban--Davenport--Halberstam (BDH) variances distinguish the
zero residue.  Translating a physical interval therefore replaces the deleted
zero residue by a modulus-dependent moving hole.  We isolate and prove the
exact correction.  Deleting a residue is an orthogonal rank-one loss from the
all-residue variance; changing the deleted residue is consequently rank two,
with exact nonzero spectrum
$\pm\sqrt{q(q-2)}/(q-1)$.  We lift the identity through the mandatory
$(q-2)/(q-1)$ diagonal subtraction and through four-packet complex
polarization.  For length-$O(H)$ bounded-overlap blocks, moduli
$q\in(Q,2Q]$, and a Schwartz localized kernel, we then prove the deterministic
bound
\[
 \sum_{b,c}|M_{bc}|\ll A_\beta A_w J(H^2+HQ+Q^2).
\]
At the frozen scales $H=x^{21/32}$, $Q=x^{1/3}$ and
$J=x/H\,x^{o(1)}$, this is $x^{53/32+o(1)}=
x^{5/3-1/96+o(1)}$.  Thus the translation subgate, and only that subgate,
pays the critical $1/96$ clock and the required strict margin for every fixed
$1/400<\delta'<1/96$.  The claim is structural L1: we do not prove the
prime-only, $q$-weighted, kernel-localized signed zero-hole BDH theorem, a
full Gate-B estimate, an arithmetic advance, an $L^2$ estimate, fixed-atom
credit, or a twin-prime theorem.
\end{abstract}

\section{Introduction and claim firewall}

The V59 prime-dynamics reduction expresses one signed endpoint scalar as a
four-packet polarization of reduced-residue BDH remainders
\cite{V59Artifact}.  A direct comparison with general-sequence BDH results is
not translation invariant: after writing a physical variable as $n=s+m$,
the condition $q\nmid n$ deletes the local residue $m\equiv-s\pmod q$, not
the local zero residue.  The purpose of this paper is to remove that
translation mismatch exactly and to determine its full deterministic cost.

The answer has two parts.  First, a finite-dimensional projector identity
shows that the variance correction is rank two.  The exact diagonal-subtracted
remainder additionally requires a two-residue energy correction.  Second,
kernel localization and block geometry, rather than finite rank, yield the
needed power: the correction reassembles with $J$ effective blocks, not
$J^2$.

The status language throughout is
\begin{align*}
 \texttt{V60\_ROUTE\_ADVANCE}&=\texttt{YES},\\
 \texttt{V60\_TRANSLATION\_SUBGATE\_DELTA}&=\texttt{1\_OVER\_96\_PROVED},\\
 \texttt{V60\_TRANSLATION\_SUBGATE\_STRICT\_1\_OVER\_400}&=\texttt{PAID},\\
 \texttt{V60\_FULL\_GATE\_B\_STRICT\_1\_OVER\_400}&=\texttt{UNPAID},\\
 \texttt{V60\_ARITHMETIC\_ADVANCE}&=\texttt{NO},\\
 \texttt{V60\_FIXED\_ATOM\_CREDIT}&=0,\\
 \texttt{V60\_L2}&=\texttt{NONE},\qquad
 \texttt{TPC\_207\_TRIGGER}=\texttt{true}.
\end{align*}
No statement below is a twin-prime theorem.  We use ``prove'' for the exact
algebra and deterministic estimates established here; source comparisons are
explicitly separated.

\section{Reduced-residue BDH rows}

Fix $q\ge2$ and a complex row $z=(z_r)_{r\in\Fq}$.  Put
\[
 \mu=\frac1q\sum_{r\in\Fq}z_r,
 \qquad
 V_{\rm all}(z)=\sum_{r\in\Fq}|z_r-\mu|^2.
\]
For a deleted residue, or \emph{hole}, $h\in\Fq$, define
\[
 \mu_h=\frac1{q-1}\sum_{r\ne h}z_r,
 \qquad
 V_h(z)=\sum_{r\ne h}|z_r-\mu_h|^2.
\]
For a prime modulus, the gcd-grouped variance used by Harper
\cite{Harper2024} has a singleton non-unit class and hence zero contribution
from that class.  Its prime row is therefore exactly $V_0$.  The translated
physical row will instead be $V_{h_q}$ for a modulus-dependent $h_q$.

We also require a diagonal ledger.  Let $E_r\ge0$ denote the coefficient
energy in residue $r$, and set
\begin{equation}\label{eq:remainder}
 \kappa_q=\frac{q-2}{q-1},
 \qquad
 R_h(z,E)=V_h(z)-\kappa_q\sum_{r\ne h}E_r.
\end{equation}
The factor $q-2$ records the exact number of nonprincipal characters after
the reduced-residue character expansion; replacing it by $q-1$ changes the
scalar.

\section{The moving-hole projector theorem}

\begin{theorem}[Moving-hole identity]\label{thm:moving-hole}
For every $q\ge2$, $z\in\mathbb C^q$, and $h\in\Fq$,
\begin{equation}\label{eq:main-identity}
 \boxed{V_h(z)=V_{\rm all}(z)-\frac q{q-1}|z_h-\mu|^2.}
\end{equation}
Consequently,
\begin{equation}\label{eq:hole-difference}
 V_h(z)-V_0(z)=\frac q{q-1}
 \left(|z_0-\mu|^2-|z_h-\mu|^2\right).
\end{equation}
\end{theorem}

\begin{proof}
Write $u_r=z_r-\mu$.  Then $\sum_r u_r=0$ and
\[
 \mu_h=\mu-\frac{u_h}{q-1}.
\]
For $r\ne h$ we have $z_r-\mu_h=u_r+u_h/(q-1)$.  Since
$\sum_{r\ne h}u_r=-u_h$,
\begin{align*}
 V_h(z)
 &=\sum_{r\ne h}\left|u_r+\frac{u_h}{q-1}\right|^2\\
 &=\sum_{r\ne h}|u_r|^2
   -\frac{2}{q-1}|u_h|^2+\frac1{q-1}|u_h|^2\\
 &=V_{\rm all}(z)-\frac q{q-1}|u_h|^2.
\end{align*}
Subtracting the cases $h$ and $0$ gives \eqref{eq:hole-difference}.
\end{proof}

Let $P=I-q^{-1}\one\one^*$ be the all-row centering projector and define
\[
 v_h=\sqrt{\frac q{q-1}}\left(e_h-\frac1q\one\right).
\]
Then $\|v_h\|=1$, $v_h\perp\one$, and Theorem
\ref{thm:moving-hole} becomes
\begin{equation}\label{eq:projector-form}
 V_h(z)=\langle z,Pz\rangle-|\langle z,v_h\rangle|^2.
\end{equation}

\section{Sharp rank-two spectrum and obstruction}

\begin{proposition}[Exact spectrum]\label{prop:spectrum}
Suppose $h\ne0$ and $q>2$.  The self-adjoint translation-defect operator
\[
 T_h=P_{v_0}-P_{v_h},\qquad P_v=vv^*,
\]
has rank two and nonzero spectrum
\begin{equation}\label{eq:spectrum}
 \boxed{\operatorname{spec}_{\ne0}(T_h)=
 \left\{\frac{\sqrt{q(q-2)}}{q-1},
 -\frac{\sqrt{q(q-2)}}{q-1}\right\}.}
\end{equation}
For $h=0$ or $q=2$, the defect is zero.
\end{proposition}

\begin{proof}
For $h\ne0$,
\[
 \langle v_0,v_h\rangle=-\frac1{q-1}.
\]
The difference of two rank-one orthogonal projectors has trace zero and is
supported on $\operatorname{span}(v_0,v_h)$.  Its two eigenvalues are
$\pm\sqrt{1-|\langle v_0,v_h\rangle|^2}$, which gives
\eqref{eq:spectrum}.  If $q=2$, the two centered unit vectors differ only by
sign, hence their projectors coincide.
\end{proof}

In particular,
\[
 \|T_h\|=\frac{\sqrt{q(q-2)}}{q-1}\longrightarrow1.
\]
Thus finite rank gives no asymptotic saving.  This is a sharp obstruction to
an argument that merely calls the correction ``low rank.''  A related
coefficient-level obstruction is obtained by placing $L$ equal coefficients
$A$ in a single residue: when the hole is moved across that residue, the
diagonal-subtracted defect can have size
\[
 \frac{q-2}{q-1}L(L-1)|A|^2.
\]
The block theorem below must therefore use residue counts, kernel decay, and
collective geometry.

\section{Exact diagonal-subtracted lift}

\begin{theorem}[Mandatory energy correction]\label{thm:diagonal}
For the remainder \eqref{eq:remainder},
\begin{equation}\label{eq:diagonal-lift}
 \boxed{
 R_h-R_0=\frac q{q-1}
 \left(|z_0-\mu|^2-|z_h-\mu|^2\right)
 +\frac{q-2}{q-1}(E_h-E_0).}
\end{equation}
If the ambient BDH row carries outer weight $q$, then every term on the
right of \eqref{eq:diagonal-lift} must also be multiplied by $q$.
\end{theorem}

\begin{proof}
The variance difference is \eqref{eq:hole-difference}.  The diagonal
difference is
\[
 -\kappa_q\sum_{r\ne h}E_r+\kappa_q\sum_{r\ne0}E_r
 =\kappa_q(E_h-E_0).
\]
Adding the two identities proves the claim.
\end{proof}

The energy correction is essential.  For $q=5$, take residue-$0$
coefficients $5,5$, residue-$1$ coefficient $1$, and no others.  Then
\[
 z=(10,1,0,0,0),\qquad E=(50,1,0,0,0),
\]
and exact calculation gives
\begin{equation}\label{eq:q5-fixture}
 R_0=0,\qquad R_1=\frac{75}{2}.
\end{equation}
By contrast, the one-coefficient realization $E=(100,1,0,0,0)$ has
$R_0=R_1=0$; the row $z$ alone cannot certify a remainder mismatch.

\section{Translation and four-packet polarization}

Fix the physical convention
\begin{equation}\label{eq:translation}
 n=s+m.
\end{equation}
Then $q\mid n$ if and only if $m\equiv-s\pmod q$.  Hence the moving hole is
\begin{equation}\label{eq:moving-hole-sign}
 \boxed{h_q=-s\pmod q.}
\end{equation}
For a local packet $a_s(m)=a(s+m)$ and frequency $v$, write
\[
 z^{a,s}_{q,r}(v)=\sum_{m\equiv r\,(q)}a_s(m)\e(vm/H).
\]
The physical row differs by the common unit phase $\e(vs/H)$, which cancels
in every variance and absolute square.  If the two polarized packets were
first written at own origins $s_\beta,s_w$, with $d=s_w-s_\beta$, then the
common-origin conversion is
\begin{equation}\label{eq:common-origin}
 W_r^{\rm common}(v)=\e(vd/H)W_{r-d}^{\rm own}(v).
\end{equation}
Both the relative residue shift and the phase are mandatory.

Let $B_r,W_r$ be the common-origin residue rows for an ordered block pair
$(b,c)$, and put
\[
 \bar B=\frac1q\sum_rB_r,\qquad
 \bar W=\frac1q\sum_rW_r,\qquad
 F_r=\sum_{m\equiv r\,(q)}\beta_b(m)\overline{w_c(m)}.
\]
Define $a^{(j)}=\beta+i^jw$, $j=0,1,2,3$.  The elementary complex
polarization identity in the required orientation is
\[
 X\overline Y=\frac14\sum_{j=0}^3i^j|X+i^jY|^2.
\]
Applying it before any triangle inequality to Theorem \ref{thm:diagonal}
gives the exact ordered-pair moving-hole defect
\begin{align}
 M_{bc}=\int\psi(v)\sum_{q\in\cQ}q\Bigg\{&
 \frac q{q-1}\Big[(B_0-\bar B)\overline{(W_0-\bar W)}
 \notag\\[-2mm]
 &\hspace{31mm}-(B_{h_q}-\bar B)
 \overline{(W_{h_q}-\bar W)}\Big]
 \notag\\
 &+\frac{q-2}{q-1}(F_{h_q}-F_0)\Bigg\}\,dv.
 \label{eq:Mbc}
\end{align}
Equivalently, \eqref{eq:Mbc} is one quarter of the signed $i^j$-sum of the
four scalar remainder defects.  It is not obtained by four absolute-value
estimates.

For later use, with
$\widehat B(k)=\sum_rB_r\e_q(-kr)$, define
$L_r=(B_r-\bar B)\overline{(W_r-\bar W)}$.  Before the prefactor
$q/(q-1)$ in \eqref{eq:Mbc}, the bare leverage difference is
\begin{equation}\label{eq:dft-selector}
 L_0-L_{h_q}=\frac1{q^2}\sum_{k,l\ne0}\widehat B(k)
 \overline{\widehat W(l)}\bigl(1-\e_q((k-l)h_q)\bigr).
\end{equation}
Consequently the complete leverage correction inside the braces of
\eqref{eq:Mbc} has coefficient $1/[q(q-1)]$ in front of the double sum.
After the outer BDH weight $q$ is included, its coefficient in $M_{bc}$ is
$1/(q-1)$.  The selector vanishes when $k=l$.  This off-equal-frequency
cancellation belongs to the leverage term; the separate diagonal term
$F_{h_q}-F_0$ does not disappear from \eqref{eq:Mbc}.

\section{Critical deterministic block bound}

We now state a self-contained deterministic form of the block estimate.  Let
$\{I_b\}_{b\in\mathcal B}$ and $\{I_c\}_{c\in\mathcal B}$ be intervals of
length at most $C_0H$, with bounded overlap and ordered centers separated on
the $H$ scale.  There are $J=|\mathcal B|$ effective blocks after discarding
a Schwartz tail.  Assume
\[
 |\beta_b(m)|\le A_\beta\one_{I_b}(m),\qquad
 |w_c(n)|\le A_w\one_{I_c}(n).
\]
The frequency integral is assumed to produce a kernel $K_H(m-n)$ satisfying,
for every fixed $A>2$,
\begin{equation}\label{eq:kernel}
 |K_H(m-n)|\le C_A(1+|b-c|)^{-A}
\end{equation}
whenever $m\in I_b,n\in I_c$.  This normalization absorbs the fixed smooth
cutoffs and is the one used in the V59 local scalar.  For the cross-diagonal
term we assume the usual bounded-overlap property: $I_b\cap I_c$ is empty
unless $|b-c|\le C_1$.

\begin{theorem}[Moving-hole block payment]\label{thm:block}
Let $2\le Q\le H$, and let $\cQ$ be any set of integer moduli in $(Q,2Q]$.
Under the preceding hypotheses, the defects \eqref{eq:Mbc} satisfy
\begin{equation}\label{eq:block-bound}
 \boxed{
 \sum_{b,c\in\mathcal B}|M_{bc}|
 \ll A_\beta A_w J(H^2+HQ+Q^2),}
\end{equation}
where the implied constant depends only on the support, overlap and Schwartz
constants, not on the locations of the holes $h_q$.
\end{theorem}

\begin{proof}
For $r\in\Fq$ introduce the centered residue selector
\[
 \lambda_{q,r}(m)=\one_{m\equiv r\,(q)}-\frac1q.
\]
On any interval of length at most $C_0H$,
\begin{equation}\label{eq:l1-selector}
 \sum_m|\lambda_{q,r}(m)|\ll \frac Hq+1,
\end{equation}
uniformly in $r$.  Indeed, the residue class occurs $O(H/q+1)$ times and the
constant part has total mass $O(H/q)$.

Integrate in $v$ first.  A leverage coordinate such as
$(B_r-\bar B)\overline{(W_r-\bar W)}$ becomes a double coefficient sum with
two centered selectors and kernel $K_H(m-n)$.  Equations
\eqref{eq:kernel} and \eqref{eq:l1-selector} therefore give, uniformly in the
selected residue $r$,
\[
 \left|\int\psi(v)(B_r-\bar B)
 \overline{(W_r-\bar W)}\,dv\right|
 \ll A_\beta A_w(1+|b-c|)^{-A}\left(\frac Hq+1\right)^2.
\]
The two leverage coordinates in \eqref{eq:Mbc} obey the same estimate.  The
outer $q$ is retained.  Summing the Schwartz separation first yields
\[
 \sum_{b,c}(1+|b-c|)^{-A}\ll J,
\]
not $J^2$.  Hence the total leverage contribution is
\begin{align*}
 &\ll A_\beta A_w J
 \sum_{q\in\cQ}q\left(\frac Hq+1\right)^2\\
 &\ll A_\beta A_w J
 \left(H^2\sum_{Q<q\le2Q}\frac1q
 +H\#\cQ+\sum_{q\in\cQ}q\right)\\
 &\ll A_\beta A_w J(H^2+HQ+Q^2).
\end{align*}
No prime number theorem is used; $\#\cQ\le Q$ suffices.

For the cross-diagonal, only $O(J)$ block pairs overlap.  In each such pair,
the residue support has size $O(H/q+1)$, so its total contribution is
\[
 \ll A_\beta A_wJ\sum_{q\in\cQ}q\left(\frac Hq+1\right)
 \ll A_\beta A_wJ(HQ+Q^2),
\]
which is already contained in \eqref{eq:block-bound}.
\end{proof}

\begin{remark}
The proof uses absolute values only after the exact four-packet polarization
has produced the complete ordered-pair defect $M_{bc}$.  It does not apply
triangle inequalities separately to the four packets or destroy their
algebraic cancellation.
\end{remark}

\section{The V59 critical scale}

At the frozen V59 scales
\[
 H=x^{21/32},\qquad Q=x^{1/3},\qquad
 J=\frac{x}{H}x^{o(1)},
\]
we have $Q\le H$ and
\begin{align*}
 JH^2&=xH\,x^{o(1)}=x^{53/32+o(1)},\\
 JHQ&=xQ\,x^{o(1)}=x^{4/3+o(1)},\\
 JQ^2&=x^{97/96+o(1)}.
\end{align*}
The first term dominates, and
\begin{equation}\label{eq:critical-clock}
 \frac{53}{32}=\frac53-\frac1{96}.
\end{equation}
Therefore Theorem \ref{thm:block} pays the moving-hole translation defect at
the exact $1/96$ clock.  The harmless $x^{o(1)}$ factor leaves the required
strict margin for every fixed
\begin{equation}\label{eq:strict-saving}
 \frac1{400}<\delta'<\frac1{96}.
\end{equation}

This payment is local to the translation correction.  It does not estimate
the standard zero-hole remainder $R_0$, and so it cannot be transferred to
the full Gate-B scalar by itself.

\section{Source boundary}

Harper defines a gcd-grouped variance for general sequences and proves BDH
asymptotics under explicit progressions, non-concentration, sparsity or
integer-resemblance hypotheses \cite{Harper2024}.  For a prime modulus, the
gcd-$q$ group is the singleton zero class and has zero variance; thus the
prime row is exactly the zero-hole variance $V_0$.  Theorem
\ref{thm:moving-hole} supplies the exact crosswalk from this row to the
translated physical row.

The source does not, however, prove the object still required by V59:
\begin{itemize}[leftmargin=2em]
 \item the modulus subset is prime-only and carries outer weight $q$;
 \item the variance is kernel-localized and has the exact $q-2$ diagonal
       removed;
 \item four literal complex packets occur in one signed polarization;
 \item the required input hypotheses must hold uniformly in block and
       frequency; and
 \item the local estimates must be reassembled collectively without an
       internal modulus, packet, or block triangle loss.
\end{itemize}
Moreover, after exact diagonal subtraction the signed remainder is not
positive, so an all-moduli estimate cannot simply be restricted to primes.
We therefore source-lock the conclusion as follows: the Harper translation
mismatch is no longer a fatal obstruction, but Harper is not a direct theorem
attachment for the full Gate-B object.

\section{Executable exact certificate}

The accompanying standard-library artifact represents complex numbers as
pairs of rational numbers and checks, without floating-point tolerances:
\begin{itemize}[leftmargin=2em]
 \item all-hole identities for $q=2,3,5,7$;
 \item the exact diagonal lift and translation sign;
 \item four-packet polarization with weights $i^j/4$;
 \item the squared spectrum $q(q-2)/(q-1)^2$ and degeneracies;
 \item the corrected fixture \eqref{eq:q5-fixture};
 \item the rational exponent identity \eqref{eq:critical-clock}; and
 \item the explicit claim firewall.
\end{itemize}
An independent checker reimplements the arithmetic without importing the
main implementation.  Both normal and optimized Python modes are required to
produce identical output.  These finite checks are quality assurance and
falsifiers for signs, factors, and fixtures; they are not evidence for the
general theorems, whose proofs appear above.

\section{Limitations and open theorem}

The strongest positive result is the exact decomposition
\begin{align*}
 \text{physical moving-hole remainder}
 ={}&\text{standard zero-hole remainder}\\
 &+\text{selected leverage defect}\\
 &+\text{selected energy defect},
\end{align*}
together with a deterministic $x^{53/32+o(1)}$ payment for the complete
translation defect.  The strongest obstruction is equally explicit: the
rank-two operator norm tends to one, so no saving is available from finite
rank alone.

The first fatal remaining theorem is:
\begin{quote}
Prove a power-saving estimate for the standard zero-hole, prime-only,
$q$-weighted, kernel-localized, exact-diagonal-subtracted signed remainder of
the four literal V59 packets, uniformly over blocks and frequencies, and
compile the block family collectively at a fixed saving greater than
$1/400$.
\end{quote}
The DFT selector $1-\e_q((k-l)h_q)$ in
\eqref{eq:dft-selector} is a reusable clue: it removes equal frequencies from
the translation leverage term and may expose a post-dispersion or
Kloosterman-cell interface.  That arithmetic step remains open.

To conclude, TPC-207 is a proved structural L1 route advance.  It proves the
exact moving-hole projector and $q-2$ diagonal compiler and pays the
translation subgate for every fixed $1/400<\delta'<1/96$.  It proves no full
Gate-B theorem, no global arithmetic advance, no $L^2$ estimate, no
fixed-atom credit, and no twin-prime theorem.

\bibliographystyle{plain}
\bibliography{references}

\end{document}
